\documentclass[11pt, a4paper]{article}
\usepackage[utf8]{inputenc}
\usepackage{amsmath, amssymb}
\usepackage[margin=1in]{geometry}
\usepackage{fancyhdr}
\usepackage{xcolor}
\definecolor{UNIBBlue}{RGB}{0, 51, 102}
\pagestyle{fancy}
\fancyhf{}
\rhead{\textbf{Optimisasi untuk Teknik Elektro}}
\lhead{Catatan Kuliah - Minggu 6}
\cfoot{\thepage}
\title{\vspace{-1cm}\color{UNIBBlue}\textbf{Konsep Konveksitas}\vspace{-0.5cm}}
\author{Ir. Novalio Daratha S.T., M.Sc., Ph.D.}
\date{Minggu 6}
\begin{document}
\maketitle
\section{Pengantar Konveksitas}
Dalam optimisasi, pemisahan yang paling penting bukan antara masalah Linier dan Non-Linier, melainkan antara masalah \textbf{Cembung (Convex)} dan \textbf{Tidak Cembung (Non-Convex)}. Optimisasi cembung menjamin bahwa setiap minimum lokal adalah minimum global.
\section{Definisi Formal}
\textbf{Himpunan Cembung:} Suatu himpunan $C$ disebut cembung jika untuk setiap $x, y \in C$ dan setiap $\theta$ dengan $0 \le \theta \le 1$, berlaku:
$\theta x + (1 - \theta)y \in C$.
\textbf{Fungsi Cembung:} Fungsi $f(x)$ disebut cembung jika:
$f(\theta x + (1-\theta)y) \le \theta f(x) + (1-\theta)f(y)$.
\end{document}